% =====================================================================
\section{Variational numerical method}\label{sec:numerical}
% =====================================================================

\subsection{Incremental minimisation}\label{sec:incremental}

We approximate energetic solutions by incremental
minimisation~\cite{ortiz1999tq,mielke2008zb,conti2008,romero2021dd}. Partition
$[0,T]$ as $0=t_0<t_1<\dots<t_N=T$, write $h_k=t_{k+1}-t_k$,
$S_k=S(t_k)$, and, given $\bs q_k$, define
\begin{equation}\label{eq:Ik}
  \bs q_{k+1}\in\argmin_{\bs q\in\Q} I_k(\bs q),
  \qquad
  I_k(\bs q)=E(\bs q,S_{k+1})-E(\bs q_k,S_k)+h_k\,\Psi\!\Bigl(\tfrac{\bs q-\bs q_k}{h_k}\Bigr).
\end{equation}
By the $1$-homogeneity of $\Psi$ and \cref{prop:Dr_linear}, minimising
$I_k$ is equivalent to minimising
\begin{equation}\label{eq:Ikstar}
  I_k^\ast(\bs q)=E(\bs q,S_{k+1})+\Dr(\bs q_k,\bs q),
\end{equation}
whose optimality condition is the backward-Euler discretisation
of~\eqref{eq:biot},
\begin{equation}\label{eq:discrete_inclusion}
  \bs 0\in\frac{\partial E}{\partial\bs q}(\bs q_{k+1},S_{k+1})
  +\partial\Psi(\bs q_{k+1}-\bs q_k).
\end{equation}

\begin{proposition}[Return map]\label{prop:return_map}
Under~\ref{A:Psi}, the incremental problem~\eqref{eq:Ikstar} is
equivalent to~\eqref{eq:discrete_inclusion}, which splits into two
mutually exclusive cases:
\begin{enumerate}[label=\textnormal{(\roman*)},leftmargin=2.4em,itemsep=2pt]
  \item \emph{Elastic lock:} $\bs q_{k+1}=\bs q_k$ iff
  $-\partial_{\bs q}E(\bs q_k,S_{k+1})\in\E$;
  \item \emph{Dissipative flow:} $\bs q_{k+1}\ne\bs q_k$ iff
  $-\partial_{\bs q}E(\bs q_{k+1},S_{k+1})\in\partial\Psi(\bs q_{k+1}-\bs q_k)\setminus\E$.
\end{enumerate}
In the scalar case $\Psi(\dot q)=\rho|\dot q|$, with trial force
$f^{\mathrm{tr}}=\ell(S_{k+1})-g(q_k)$, the update reads
\begin{equation}\label{eq:scalar_return_map}
  q_{k+1}=
  \begin{cases}
    q_k, & |f^{\mathrm{tr}}|\le\rho \quad(\text{lock}),\\[2pt]
    \text{solution of } g(q_{k+1})=\ell(S_{k+1})-\sgn(f^{\mathrm{tr}})\,\rho,
      & |f^{\mathrm{tr}}|>\rho\quad(\text{flow}).
  \end{cases}
\end{equation}
\end{proposition}

Equation~\eqref{eq:scalar_return_map} is a two-sided return
map~\cite{Simo1986,hanreddy2013}: a trial (elastic) predictor followed,
if the yield surface $|f|=\rho$ is violated, by a plastic corrector
onto it.

\subsection{Discrete stability and energy inequality}\label{sec:discrete_stab}

\begin{proposition}[Discrete stability and energy
bound]\label{prop:discrete_energy}
Let $\{\bs q_k\}$ be generated by~\eqref{eq:Ikstar}. Then for each $k$,
\begin{align}
  &E(\bs q_{k+1},S_{k+1})\le E(\tilde{\bs q},S_{k+1})+\Dr(\bs q_k,\tilde{\bs q})
  \quad\forall\,\tilde{\bs q}\in\Q,
  \label{eq:disc_stab}\\
  &E(\bs q_{k+1},S_{k+1})+\Dr(\bs q_k,\bs q_{k+1})
  \le E(\bs q_k,S_{k+1})
  = E(\bs q_k,S_k)+\bigl[E(\bs q_k,S_{k+1})-E(\bs q_k,S_k)\bigr].
  \label{eq:disc_energy}
\end{align}
\end{proposition}

The bound~\eqref{eq:disc_energy} is generally a strict inequality: the
scheme may over-dissipate at the discrete level because the external
work is evaluated at the frozen state $\bs q_k$. As $h\to0$ the gap
telescopes to the exact continuous balance~\eqref{eq:energy_balance},
which is the content of the next result.

\subsection{Convergence}\label{sec:convergence}

\begin{theorem}[Convergence of the incremental scheme]\label{thm:convergence}
Under~\ref{A:W}--\ref{A:Psi}, let $\bs q_0\in\mathcal S(0)$ and let
$\bs q^h$ denote the piecewise-constant interpolant of the sequence
$\{\bs q_k\}$ generated by~\eqref{eq:Ikstar} on a partition of maximal
step $h$. Then, as $h\to0$, there is a subsequence $h_j\to0$ such that
$\bs q^{h_j}(t)\to\bs q(t)$ for every $t\in[0,T]$, where $\bs q$ is an
energetic solution of $(\Q,E,\Psi)$. If, moreover, the solution is
unique (\cref{thm:uniqueness} or~\cref{thm:BV_uniqueness}), the whole
family $\bs q^h$ converges.
\end{theorem}

\subsection{Consistency and error bounds}\label{sec:consistency}

Substituting the exact solution into the discrete energy
bound~\eqref{eq:disc_energy} produces, since the exact solution obeys
the continuous balance~\eqref{eq:energy_balance}, the \emph{local
truncation error}
\begin{equation}\label{eq:truncation}
  \tau_k := \int_{t_k}^{t_{k+1}}\!\partial_S E(\bs q(s),S(s))\,\dot S(s)\dd s
  -\bigl[E(\bs q(t_k),S_{k+1})-E(\bs q(t_k),S_k)\bigr],
\end{equation}
the discrepancy between the continuous work and its frozen-state
discrete approximation.

\begin{proposition}[Local consistency]\label{prop:consistency}
Assume~\ref{A:ell}, $S\in W^{1,\infty}(0,T)$, and let $\bs q$ be the
exact solution.
\begin{enumerate}[label=\textnormal{(\roman*)},leftmargin=2.4em,itemsep=2pt]
  \item \emph{Weak regularity.} If $\bs q\in\mathrm{AC}(0,T;\Q)$, then
  \begin{equation}\label{eq:cons_weak}
    |\tau_k|\le C_{\ell'}\|\dot S\|_{L^\infty}
    \Bigl(\int_{t_k}^{t_{k+1}}\|\dot{\bs q}\|\dd\tau\Bigr)h_k
    = O(h_k).
  \end{equation}
  \item \emph{Strong regularity.} If $\bs q\in W^{1,\infty}(0,T;\Q)$,
  then
  \begin{equation}\label{eq:cons_strong}
    |\tau_k|\le\tfrac12\,C_{\ell'}\|\dot{\bs q}\|_{L^\infty}\|\dot S\|_{L^\infty}\,h_k^2
    = O(h_k^2).
  \end{equation}
\end{enumerate}
\end{proposition}

\begin{theorem}[Global energy consistency]\label{thm:global_consistency}
Let $h=\max_k h_k$ and
$\mathcal E_\tau:=\sum_{k=0}^{N-1}|\tau_k|$. Under the hypotheses of
\cref{prop:consistency},
\begin{equation}\label{eq:global_cons}
  \mathcal E_\tau\le C\,h,\qquad
  C=\begin{cases}
    C_{\ell'}\|\dot S\|_{L^\infty}\int_0^T\|\dot{\bs q}\|\dd\tau, & \bs q\in\mathrm{AC},\\[4pt]
    \tfrac12\,C_{\ell'}\|\dot{\bs q}\|_{L^\infty}\|\dot S\|_{L^\infty}\,T, & \bs q\in W^{1,\infty},
  \end{cases}
\end{equation}
so the scheme is globally first-order consistent in the energy balance,
irrespective of the regularity of $\bs q$.
\end{theorem}

The regularity of $\bs q$ is inherited from the smoothness of the data
$W,\bs\ell,S$ away from jumps; at a jump $\bs q$ is only of bounded
variation and the local order degrades, but the global energy
consistency $\mathcal E_\tau=O(h)$ persists. All proofs of this section
are given in \cref{app:numerical}.
